\subsection{Lettura congiunta dei risultati}

Le campagne A--F1 mostrano perché la qualità di una mesh adattiva non possa essere valutata soltanto sulla riduzione di un funzionale geometrico o d'interpolazione. Stage A migliora l'errore d'interpolazione sui campi gaussiani, ma tale miglioramento non si trasferisce automaticamente all'errore PDE. Stage D, con target trajectory-aware, riduce invece l'errore relativo $L^2$ integrato e finale rispetto alla mesh fissa nei 28 casi esaminati. F1 estende il confronto a tre controllori classici di $r$-refinement e registra un ulteriore vantaggio medio di GDM sull'endpoint primario.

Questa evidenza è coerente con l'ipotesi che l'informazione temporale nei target sia utile per una struttura mobile. Non dimostra tuttavia che la sola architettura GNN sia responsabile del guadagno: il risultato appartiene alla pipeline composta da feature dello stato FEM corrente, rete, filtro geometrico, interpolazione temporale dello spostamento e solver ALE. Inoltre, i target Stage D sono riferimenti a budget limitato costruiti con la MMS; non sono minimizzatori convergenti e non costituiscono un oracle certificato.

\subsection{Posizionamento rispetto alla letteratura}

La combinazione GNN e movimento nodale è già presente in M2N, DMM, UM2N, G-Adaptivity e UGM2N \citep{song2022m2n,hu2024dmm,zhang2024towards,rowbottom2025gadaptivity,wang2025ugm2n}. Non rivendichiamo priorità per il mesh movement neurale, per la topologia fissa, per l'ALE o per la prevenzione del tangling, considerati isolatamente. Il contributo da valutare è l'integrazione verificabile di un mesh mover supervisionato su finestre temporali, di un filtro geometrico pathwise e di un solver ALE con diagnostiche D-GCL, insieme al confronto con strategie classiche a parità di gradi di libertà nel problema studiato. La novità assoluta di questa combinazione richiede comunque una revisione sistematica della letteratura più ampia di quella qui riassunta.

\subsection{Accuratezza, geometria e safety filter}

Il principale compromesso emerso in F1 riguarda la geometria. GDM ottiene errori $L^2$ inferiori nel test considerato, ma genera elementi con qualità minima più bassa e aspect ratio più elevati rispetto ai controllori classici. Il suo minimo Jacobiano si avvicina alla soglia assoluta configurata. La diagnostica validation evidenzia che il filtro interviene nell'80\% degli eventi osservati e conserva circa il 48.1\% della norma aggregata degli spostamenti proposti. Una parte delle proposte raw contiene triangoli invertiti, ma il filtro impedisce che tali proposte vengano applicate integralmente.

Questo comportamento può essere interpretato come una dipendenza non trascurabile del controller da un meccanismo esterno di ammissibilità. Non è un difetto da nascondere né una prova che il filtro sia scorretto: è una proprietà della pipeline che va caratterizzata. In particolare, la soglia $J_e\ge10^{-4}$ è assoluta e dipende dalla scala dell'area elementare; non è stata dimostrata indipendente dalla risoluzione. Prima di estendere il metodo a mesh più fini occorre formulare e verificare un criterio geometrico scalato o relativo, e dimostrare la massimalità del fattore $\alpha$ con test su mesh distorte e casi quasi degeneri. La finalizzazione F1 ha inoltre evidenziato una discrepanza di arrotondamento inferiore a $3\times10^{-19}$ sotto la soglia; il gate numerico deve essere reso coerente con la tolleranza esplicita del core, senza modificare la soglia per migliorare artificialmente i risultati.

La diagnostica di tracking mostra un arricchimento locale mediano di circa il 21.3\% in quattro casi validation. È un'indicazione che la mesh redistribuisce effettivamente i nodi verso la regione mobile, ma non dimostra tracking perfetto né viene usata come evidenza held-out indipendente. La distanza tra il centroide della concentrazione di elementi e il centro gaussiano rimane non nulla e cresce nel corso della traiettoria; perciò nel manoscritto si parla di adattamento dinamico della risoluzione, non di inseguimento esatto del massimo.

\subsection{Verifica numerica, SUPG e D-GCL}

E1 documenta un ordine temporale che si avvicina a uno e ordini spaziali compatibili con $P_1$ sulla coppia $16^2$--$32^2$. Il livello $64^2$ è più accurato in valore assoluto, ma il rapporto operativo di contaminazione temporale supera il limite del 10\%; l'ordine tra $32^2$ e $64^2$ non viene pertanto considerato confermato. La curva work--precision a $\Delta t$ fisso non va interpretata come separazione pura dell'errore spaziale. La SUPG è stata implementata come opzione nel core, ma nella configurazione MMS testata non ha ridotto l'errore rispetto a Galerkin; questa evidenza non può essere generalizzata a ogni regime di Péclet. La D-GCL è verificata in un test indipendente con residui geometrici e algebrici prossimi alla precisione macchina e con free-stream compatibile con la tolleranza effettiva del solve. Tale verifica non coincide con una dimostrazione del bilancio integrale completo della PDE forzata.

\subsection{Baseline classiche e costo}

Le baseline MMPDE sono state selezionate su validation e valutate sugli stessi casi held-out. Tuttavia, 26 eventi su 840 hanno raggiunto il limite di 200 iterazioni pseudo-temporali. L'energia è monotona e non vi sono fallimenti di line-search, ma il budget iterativo può aver limitato l'ottimizzazione in alcuni eventi. La comparazione F1 va quindi considerata preliminare fino alla verifica della stabilità dei risultati aumentando il budget sulla validation e congelando poi i parametri prima di una nuova valutazione indipendente.

Il costo online mediano GDM è circa 12.6 volte quello della mesh fissa. GDM è più rapido degli attuali controller NumPy/Python, ma questo dato è dipendente dall'implementazione e non stabilisce una complessità inferiore rispetto ai metodi classici. Il costo offline di training non è incluso. Per una valutazione di efficienza più forte occorrono controller classici ottimizzati o nativi, misure su più risoluzioni e una contabilità del costo ammortizzato per numero di simulazioni.

\subsection{Validità esterna e ulteriori limiti}

La validità esterna è limitata a un dominio quadrato bidimensionale, una triangolazione iniziale strutturata, elementi $P_1$, coefficienti fisici fissi e una famiglia MMS gaussiana. Il test estrapolativo varia principalmente $\sigma$; non dimostra trasferimento a una PDE differente, a geometrie o condizioni al contorno non viste, né a topologie o risoluzioni diverse. Gli esperimenti Stage B--C sono diagnostiche preliminari e il confronto Stage C usa il prototipo Python/SciPy; non devono essere presentati come prove di prestazione del core nativo. La riga a tre iterazioni L-BFGS è una baseline a basso budget, non un oracle.

Il test set F1 contiene 28 casi, con confronti appaiati e correzione Holm. La significatività statistica all'interno di questo protocollo non equivale a una garanzia di generalizzazione a una popolazione più ampia. Inoltre la diagnostica di tracking è stata condotta su quattro validation case e non deve essere confusa con il test held-out.


\subsection{Riproducibilità della pipeline e artefatti}
Gli esperimenti riportano snapshot di configurazione, identificativi dei checkpoint e, nei protocolli più recenti, hash dei sorgenti e dei casi. Ciò non equivale ancora a una dimostrazione di riproducibilità da una checkout indipendente. L'audit della repository ha individuato in versioni precedenti percorsi predefiniti non uniformi per configurazioni e checkpoint, oltre a un fallback del caricamento dati che deve essere verificato in modalità fail-closed. Prima della sottomissione occorre quindi eseguire un controllo end-to-end da un ambiente pulito, verificare disgiunzione e completezza del manifest degli split, archiviare le versioni effettive di Python, PyTorch, deal.II e del compilatore e rigenerare tabelle e figure a partire dagli output grezzi. Nel manoscritto non viene pertanto rivendicata, allo stato attuale, la riproducibilità indipendente dell'intera pipeline con un singolo comando.

\subsection{Priorità prima della sottomissione}

Le attività che possono modificare sostanzialmente la solidità delle conclusioni sono: (i) chiudere la separazione spazio--tempo al livello $64^2$ e rigenerare il work--precision con errore temporale controllato; (ii) completare l'audit matematico/floating-point del safety filter e definire una soglia scalata per eventuali risoluzioni ulteriori; (iii) verificare la convergenza delle baseline MMPDE con budget maggiori, senza usare i test per la selezione; (iv) valutare condizionamento FEM e qualità geometrica; (v) rendere riproducibile la pipeline da una checkout pulita con config, checkpoint, hash, versione deal.II, output grezzi e comandi documentati. Ulteriori architetture GNN non sono necessarie prima di aver chiuso questi punti.
